\documentclass[11pt]{article}
\usepackage[margin=1.1in]{geometry}
\usepackage[T1]{fontenc}
\usepackage{lmodern}
\usepackage{amsmath,amssymb,amsthm,mathtools}
\usepackage{microtype}
\usepackage{enumitem}
\usepackage[colorlinks=true,linkcolor=black,citecolor=black,urlcolor=black]{hyperref}

\newtheorem{maintheorem}{Theorem}
\renewcommand{\themaintheorem}{\Alph{maintheorem}}
\newtheorem{lemma}{Lemma}
\theoremstyle{definition}
\newtheorem{definition}{Definition}
\theoremstyle{remark}
\newtheorem{remark}{Remark}

\newcommand{\R}{\mathbb{R}}
\newcommand{\E}{\mathbb{E}}
\newcommand{\cL}{\mathcal{L}}
\newcommand{\cB}{\mathcal{B}}
\newcommand{\cK}{\mathcal{K}}
\newcommand{\cR}{\mathcal{R}}
\newcommand{\eps}{\varepsilon}
\newcommand{\sm}{\setminus}
\newcommand{\ip}[2]{\langle #1,#2\rangle}
\newcommand{\sgn}{\operatorname{sgn}}
\newcommand{\ND}[1]{\textup{(ND#1)}}

\title{Saddle-to-saddle dynamics of deep polynomial ResNets\\[4pt]\large The case of one neuron per layer}
\author{}
\date{November 2026}

\begin{document}
\maketitle

\noindent
This note proves Theorem~A of \texttt{saddle\_to\_saddle.tex} for ResNets of width $d=1$ without biases, where every neuron has a single weight. The proof is complete, and it is shorter than the general one for two reasons. First, the direction of a neuron is its sign, which the flow preserves. Hence the dormant flow is solved in closed form, condition \ND2 is void, and the limit process is a race between clocks that run at piecewise constant speeds. Second, in the escape from a saddle there is no direction to align. Only the active weights have to be controlled, and they are slaved to the escaping weight: compared at equal values of this weight, two trajectories do not move apart. So the weighted norm and the contraction rate of the general proof are not needed.

\section{Setting}

\paragraph{Network.} Let $\sigma$ be a polynomial with $\sigma(0)=\dots=\sigma^{(k-1)}(0)=0\ne\sigma^{(k)}(0)$ for some $k\ge3$, e.g.\ $\sigma(z)=z^k$, and let
\[
h_0=x,\qquad h_n=h_{n-1}+\sigma(w_nh_{n-1})\quad(1\le n\le L),\qquad f_\theta=h_L,
\]
with $x\in\R$ and $\theta=(w_1,\dots,w_L)\in\Theta=\R^L$. Layer $n$ is a single \emph{neuron}, with weight $w_n$, and $f_\theta$ is a polynomial in $x$ of degree at most $r^L$, where $r=\deg\sigma$. Let $F(\theta)=\cL(f_\theta)$, where $\cL$ is a $C^\infty$ loss on the space $V$ of polynomials of degree at most $r^L$, e.g.\ $\cL(f)=\frac12\E|f(x)-y|^2$. Let $\theta^\alpha$ be the gradient flow $\dot\theta=-\nabla F(\theta)$ with $\theta^\alpha(0)=\alpha\xi$, where $\xi\in\Theta$ is fixed and $\alpha\to0$.

\paragraph{Notation.} For a set $A\subseteq\{1,\dots,L\}$ of (active) neurons, $S_A=\{\theta:w_n=0\ \forall n\notin A\}$, $\pi_A$ is the orthogonal projection onto $S_A$, $F_A=F|_{S_A}$, $\nabla_A=\pi_A\nabla$, and $\zeta_A(\theta)=\sum_{n\notin A}|w_n|^k$. We write $e_n$ for the $n$-th basis vector, $\partial_n$ for $\partial/\partial w_n$, and $\theta^{\sm n}=\theta-w_ne_n$. The \emph{utility} of neuron $n$ at $\theta$ is the monomial
\[
U_n(\theta;w)=c_n(\theta)\,w^k,\qquad c_n(\theta)=-\tfrac1{k!}\,\partial_n^kF(\theta^{\sm n}),\qquad w\in\R .
\]

\paragraph{Flatness.} Besides smoothness, the proof uses only
\begin{equation}\tag{H}\label{eq:H}
\partial_n^iF=0\ \text{ on }\{w_n=0\}\qquad\text{for every neuron $n$ and every $1\le i\le k-1$.}
\end{equation}
The network satisfies it: $f_\theta=\Phi_n\big(h_{n-1}+\sigma(w_nh_{n-1})\big)$, where $\Phi_n$ is the map computed by the layers after $n$, and neither $h_{n-1}$ nor $\Phi_n$ depends on $w_n$. Since $\sigma(t)=t^k\tau(t)$ for a polynomial $\tau$, we get $f_\theta-f_{\theta^{\sm n}}=w_n^k\Delta_\theta$ with $\Delta_\theta\in V$ polynomial in $\theta$. Hence $F(\theta)-F(\theta^{\sm n})=w_n^k\,g(\theta)$ with $g$ smooth. Depth plays no further role: everything below holds for every $F\in C^\infty(\R^L)$ satisfying~\eqref{eq:H}. For the mean squared error, the chain rule gives, with $c_\sigma=\sigma^{(k)}(0)/k!$ and the right-hand side computed at $\theta^{\sm n}$,
\[
c_n(\theta)=c_\sigma\,\E\big[h_{n-1}(x)^k\,\Phi_n'(h_n(x))\,(y-f_\theta(x))\big].
\]
This is the correlation between the $k$-th power of the input of the neuron and the back-propagated residual. At $\theta=0$ it equals $c_\sigma\E[x^k(y-x)]$ for every $n$.

\begin{lemma}[Flatness]\label{lem:flat}
Assume \eqref{eq:H}.
\begin{enumerate}[label=\textup{(\alph*)},nosep]
\item For every neuron $n$ there is $G_n\in C^\infty(\Theta)$ with $\partial_nF=-w_n^{k-1}G_n$, and $G_n=kc_n$ on $\{w_n=0\}$.
\item On every bounded subset of $\Theta$ there is $C$ such that $|\nabla_AF(\theta)-\nabla_AF(\pi_A\theta)|\le C\zeta_A(\theta)$ for every $A$.
\end{enumerate}
\end{lemma}

\begin{proof}
Let $g\in C^\infty(\Theta)$ with $\partial_n^ig=0$ on $\{w_n=0\}$ for all $i<j$. By Taylor's formula with integral remainder, $g=w_n^j\tilde g$ with $\tilde g$ smooth and $\tilde g=\partial_n^jg/j!$ on $\{w_n=0\}$. For (a), take $g=\partial_nF$ and $j=k-1$. For (b), let $n\notin A$ and $n'\in A$. Since $\partial_{n'}$ is tangent to $\{w_n=0\}$, the derivatives $\partial_n^i\partial_{n'}F=\partial_{n'}\partial_n^iF$, $1\le i\le k-1$, vanish on this hyperplane. So $\partial_{n'}F(\theta)-\partial_{n'}F(\theta^{\sm n})$ is $w_n^k$ times a smooth function: setting $w_n$ to zero changes $\nabla_AF$ by at most $C|w_n|^k$ on bounded sets. Setting the weights $w_n$, $n\notin A$, to zero one at a time gives (b).
\end{proof}

\paragraph{Consequences.} By (a), the gradient flow reads $\dot w_n=w_n^{k-1}G_n(\theta)$. Hence every hyperplane $\{w_n=0\}$ is invariant, and \emph{the sign of each weight is constant along every trajectory}. Every $S_A$ is invariant, the flow on $S_A$ is the gradient flow of $F_A$, and critical points of $F_A$ are critical points of $F$. If $w_n=0$, then $F(\theta+we_n)=F(\theta)-U_n(\theta;w)+O(|w|^{k+1})$. Hence a critical point $\theta$ of $F_A$ with $U_n(\theta;u)>0$ for some $n\notin A$ and some sign $u\in\{\pm1\}$ is a saddle point of $F$.

\section{The limit process and the result}

A \emph{non-degenerate local minimum} of $F_A$ is a point $\theta^*\in S_A$ with $\nabla F_A(\theta^*)=0$ and $\nabla^2F_A(\theta^*)\succ0$; for $A=\emptyset$, $\theta^*=0$. Let $\theta^*$ be such a point, let $n\notin A$, and let $u\in\{\pm1\}$ with $U_n(\theta^*;u)>0$. The \emph{escape orbit} of $\theta^*$ along $(n,u)$ is the maximal gradient-flow trajectory $\gamma$ in $S_{A\cup\{n\}}$ with $u\gamma_n>0$ and $\gamma(t)\to\theta^*$ as $t\to-\infty$. It exists and is unique up to time translation (Lemma~\ref{lem:escape}(a)).

\begin{definition}[Saddle-to-saddle process]\label{def:process}
Let $u_n=\sgn\xi_n$. Set $s_0=0$, $A_0=\emptyset$, $\theta^{(0)}=0$ and $\Xi_n(0)=\xi_n$. At stage $m$, we are given $s_m$, $A_m$, a non-degenerate local minimum $\theta^{(m)}$ of $F_{A_m}$, and $\Xi_n(s_m)$ for $n\notin A_m$. Let $U^{(m)}_n=U_n(\theta^{(m)};\cdot)$. Each \emph{dormant} $\Xi_n$, $n\notin A_m$, solves
\begin{equation}\label{eq:dormant}
\frac{d}{ds}\Xi_n=k\,c_n(\theta^{(m)})\,\Xi_n^{k-1},\qquad\text{that is,}\qquad \frac{d}{ds}|\Xi_n|^{2-k}=-k(k-2)\,U^{(m)}_n(u_n)\quad\text{if }\xi_n\ne0,
\end{equation}
up to its blow-up time $s_m+T^{(m)}_n$, where
\[
T^{(m)}_n=\frac{|\Xi_n(s_m)|^{2-k}}{k(k-2)\,U^{(m)}_n(u_n)}\quad\text{if }U^{(m)}_n(u_n)>0,\qquad T^{(m)}_n=\infty\quad\text{otherwise.}
\]
Let $s_{m+1}=s_m+\min_{n\notin A_m}T^{(m)}_n$. Stage $m$ is \emph{non-degenerate} if $s_{m+1}<\infty$ and
\begin{enumerate}[label=\textup{(ND\arabic*)},leftmargin=4em,nosep]
\item the minimum is attained at a single neuron $n_{m+1}$;
\setcounter{enumi}{2}
\item the escape orbit $\gamma_{m+1}$ of $\theta^{(m)}$ along $(n_{m+1},u_{n_{m+1}})$ is defined for all $t\ge0$ and converges, as $t\to\infty$, to a non-degenerate local minimum $\theta^{(m+1)}$ of $F_{A_{m+1}}$, where $A_{m+1}=A_m\cup\{n_{m+1}\}$.
\end{enumerate}
In that case, set $\Xi_n(s_{m+1})=\lim_{s\uparrow s_{m+1}}\Xi_n(s)$ for $n\notin A_{m+1}$ and continue. Otherwise stop. The \emph{limit path} is $\bar\theta(s)=\theta^{(m)}$ for $s\in[s_m,s_{m+1})$, and $\bar s_n\in(0,\infty]$ is the time up to which $\Xi_n$ is defined, so that $\bar s_n=s_m$ if $n=n_m$.
\end{definition}

The numbering is that of the full note. Its condition \ND2, on the limit direction of the escaping neuron, is void: here the direction is the constant sign $u_{n_{m+1}}$, and $U^{(m)}_{n_{m+1}}(u_{n_{m+1}})>0$ because $T^{(m)}_{n_{m+1}}<\infty$.

\begin{remark}[A race]\label{rem:race}
Summing \eqref{eq:dormant} over the stages gives, for $s<\bar s_n$,
\[
|\Xi_n(s)|^{2-k}=|\xi_n|^{2-k}-k(k-2)\int_0^sU_n\big(\bar\theta(s');u_n\big)\,ds' .
\]
So neuron $n$ carries a clock, started at $|\xi_n|^{2-k}$, which runs down at the speed $k(k-2)\,U_n(\bar\theta;u_n)$, constant on each plateau and possibly negative. The neuron activates when its clock reaches $0$, i.e., when the utility accumulated along its sign reaches $|\xi_n|^{2-k}/(k(k-2))$. For odd $k$, a clock runs down only while $u_n$ is the sign of $c_n(\bar\theta)$. For even $k$, the sign plays no role. For instance, for the mean squared error all neurons have the same utility $c\,w^k$ at $\theta=0$, with $c=c_\sigma\E[x^k(y-x)]$. The first neuron to activate is the one with the largest $|\xi_n|$ among those with $c\,u_n^k>0$, at the time $s_1=|\xi_{n_1}|^{2-k}/(k(k-2)|c|)$, and \ND1 asks that it be unique. Its layer is decided by the initialization.
\end{remark}

\begin{maintheorem}[width one]\label{thm:main}
Let $k\ge3$ and $\xi\in\Theta$. Assume that stages $0,\dots,M-1$ are non-degenerate, and let $S\in(s_M,s_{M+1})$. For $\alpha$ small, $\theta^\alpha$ is defined on $[0,\alpha^{2-k}S]$. As $\alpha\to0$, in the time $s=\alpha^{k-2}t$:
\begin{enumerate}[label=\textup{(\roman*)},nosep]
\item \emph{(plateaus)} for every $\delta>0$, $\theta^\alpha(\alpha^{2-k}s)\to\bar\theta(s)$ uniformly on $[0,S]\sm\bigcup_{1\le m\le M}(s_m-\delta,s_m+\delta)$. In particular the loss converges to the step function $F(\bar\theta(s))$, and $F(\theta^{(0)})>\dots>F(\theta^{(M)})$;
\item \emph{(dormant neurons)} for every neuron $n$ and every $\delta>0$, the rescaled weights $\alpha^{-1}w^\alpha_n(\alpha^{2-k}s)$ converge to $\Xi_n(s)$ uniformly on $[0,\min(S,\bar s_n-\delta)]$;
\item \emph{(jumps)} for $1\le m\le M$, there are times $t^\alpha_m$ with $\alpha^{k-2}t^\alpha_m\to s_m$ such that $\theta^\alpha(t^\alpha_m+\cdot)$ converges to a time translate of $\gamma_m$, locally uniformly on $\R$. Moreover, for $\delta$ smaller than half the gaps between $0,s_1,\dots,s_M,S$, the curve $\{\theta^\alpha(\alpha^{2-k}s):|s-s_m|\le\delta\}$ converges in Hausdorff distance to $\gamma_m(\R)\cup\{\theta^{(m-1)},\theta^{(m)}\}$.
\end{enumerate}
Each $\theta^{(m)}$ is a critical point of $F$, and a saddle point for $m<M$.
\end{maintheorem}

\section{Proof}

Assume \eqref{eq:H} with $k\ge3$. In Lemmas~\ref{lem:active}--\ref{lem:escape}, $\theta^*$ is a non-degenerate local minimum of $F_A$, and $\mu,r_*>0$ are such that $\nabla^2F_A\succeq\mu I$ on $\{a\in S_A:|a-\theta^*|\le r_*\}$. Since $\nabla F_A(\theta^*)=0$, it follows that
\begin{equation}\label{eq:monotone}
\ip{a-\theta^*}{\nabla F_A(a)}\ge\mu|a-\theta^*|^2\qquad\text{for }a\in S_A,\ |a-\theta^*|\le r_* .
\end{equation}

\begin{lemma}[Active neurons]\label{lem:active}
There are $r_0\in(0,1)$ and $C_0$ such that the following holds. Let $\theta(t)$, $t\in[t_0,t_1]$, be a gradient-flow trajectory, and let $a=\pi_A\theta$. If $|a(t_0)-\theta^*|\le r_0$ and $\zeta_A(\theta)\le r_0$ on $[t_0,t_1]$, then
\[
|a(t)-\theta^*|\le e^{-\mu(t-t_0)}|a(t_0)-\theta^*|+C_0\sup_{[t_0,t]}\zeta_A(\theta)\qquad\text{for }t\in[t_0,t_1].
\]
\end{lemma}

\begin{proof}
Let $C$ be the constant of Lemma~\ref{lem:flat}(b) on the ball of radius $|\theta^*|+r_*+\sqrt L$. As long as $|a-\theta^*|\le r_*$, the point $\theta$ lies in this ball, because $|w_n|\le1$ for $n\notin A$. Then, by Lemma~\ref{lem:flat}(b) and \eqref{eq:monotone},
\[
\frac{d}{dt}\frac{|a-\theta^*|^2}2=-\ip{a-\theta^*}{\nabla_AF(\theta)}\le-\mu|a-\theta^*|^2+C\zeta_A(\theta)\,|a-\theta^*| .
\]
Grönwall's lemma gives the claim with $C_0=C/\mu$, as long as $|a-\theta^*|\le r_*$. If $(1+C_0)r_0<r_*$, the bound itself shows that $|a-\theta^*|$ never reaches $r_*$.
\end{proof}

\begin{lemma}[Dormant neurons]\label{lem:dormant}
Let $\Xi_n\colon[s_-,s_+]\to\R$, $n\notin A$, be solutions of $\Xi_n'=k\,c_n(\theta^*)\,\Xi_n^{k-1}$ with $|\Xi_n|\le R$. For small $\alpha>0$, let $\theta^\alpha$ be gradient-flow trajectories and $\tau_\alpha$ times such that $|\pi_A\theta^\alpha(\tau_\alpha)-\theta^*|\le r_0/2$ and $\alpha^{-1}w^\alpha_n(\tau_\alpha)\to\Xi_n(s_-)$ for all $n\notin A$. Write $t(s)=\tau_\alpha+\alpha^{2-k}(s-s_-)$. Then, for $\alpha$ small, $\theta^\alpha$ is defined on $[\tau_\alpha,t(s_+)]$, and:
\begin{enumerate}[label=\textup{(\alph*)},nosep]
\item $\alpha^{-1}w^\alpha_n(t(s))\to\Xi_n(s)$ uniformly on $[s_-,s_+]$ as $\alpha\to0$, for every $n\notin A$;
\item $|\pi_A\theta^\alpha(t)-\theta^*|\le e^{-\mu(t-\tau_\alpha)}r_0/2+C_0L(R+1)^k\alpha^k$ for $t\in[\tau_\alpha,t(s_+)]$.
\end{enumerate}
\end{lemma}

\begin{proof}
By Lemma~\ref{lem:flat}(a), the rescaled weights $\tilde\Xi_n(s)=\alpha^{-1}w^\alpha_n(t(s))$ solve $\tilde\Xi_n'=G_n(\theta^\alpha)\,\tilde\Xi_n^{k-1}$, with no error term. Let $s^*\le s_+$ be maximal such that, on $[s_-,s^*]$, $\theta^\alpha$ is defined and $|\tilde\Xi_n|\le R+1$ for all $n\notin A$. There, $\zeta_A(\theta^\alpha)\le L(R+1)^k\alpha^k\le r_0$ for small $\alpha$. So Lemma~\ref{lem:active} gives (b), and $\theta^\alpha$ stays in a fixed ball, on which the $G_n$ are bounded and Lipschitz. Since $G_n(\theta^*)=kc_n(\theta^*)$ for $n\notin A$,
\[
|G_n(\theta^\alpha)-kc_n(\theta^*)|\le C|\theta^\alpha-\theta^*|\le C'\big(|\pi_A\theta^\alpha-\theta^*|+\alpha\big),
\]
and by (b) the integral of the right-hand side over $s\in[s_-,s^*]$ is $O(\alpha^{k-2}+\alpha)$. Grönwall's lemma gives $|\tilde\Xi_n-\Xi_n|\le C''\big(|\tilde\Xi_n(s_-)-\Xi_n(s_-)|+\alpha^{k-2}+\alpha\big)$ on $[s_-,s^*]$. For small $\alpha$ the right-hand side is less than $1$, so $|\tilde\Xi_n|<R+1$ and the trajectory stays bounded. By maximality, $s^*=s_+$.
\end{proof}

The escape from a saddle is the heart of the proof. Let $n\notin A$, $A'=A\cup\{n\}$, and let $u\in\{\pm1\}$ with $\kappa_0:=\frac k2\,U_n(\theta^*;u)>0$. We write $\theta=a+we_n+z$ with $a=\pi_A\theta$, $w=w_n$ and $z=\theta-\pi_{A'}\theta$, which collects the other inactive neurons, and $\rho=|w|$. Let $r_0$ and $C_0$ be as in Lemma~\ref{lem:active}.

\begin{lemma}[Escape]\label{lem:escape}
For every $K\ge0$ there are $\rho_1,\eps_1\in(0,1]$ and $C$ with the following properties.
\begin{enumerate}[label=\textup{(\alph*)},nosep]
\item There is a gradient-flow trajectory $\gamma\colon(-\infty,0]\to S_{A'}$ with $u\gamma_n>0$ and $|\gamma_n(0)|=\rho_1$, such that $|\gamma_n|$ is increasing, $|\pi_A\gamma-\theta^*|\le C_0|\gamma_n|^k$, and $\gamma(t)\to\theta^*$ as $t\to-\infty$. Every gradient-flow trajectory $\tilde\gamma\colon(-\infty,t']\to S_{A'}$ with $u\tilde\gamma_n>0$ and $\tilde\gamma(t)\to\theta^*$ as $t\to-\infty$ is a time translate of $\gamma$.
\item Let $\theta(t)$ be a gradient-flow trajectory and $t_0$ a time such that $uw(t_0)>0$, $\rho_0:=\rho(t_0)<\rho_1$, $\eta:=|z(t_0)|\le\eps_1\rho_0$ and $|a(t_0)-\theta^*|\le K\rho_0^k$. Then $\rho$ increases and reaches $\rho_1$ at a time $t_1\le t_0+\rho_0^{2-k}/((k-2)\kappa_0)$. Moreover, for $t\in[t_0,t_1]$,
\[
|z(t)|\le2\eta\qquad\text{and}\qquad|\theta(t)-\gamma(t')|\le C(\rho_0^k+\eta),
\]
where $t'\le0$ is the time at which $|\gamma_n(t')|=\rho(t)$.
\end{enumerate}
\end{lemma}

\begin{proof}
\emph{Step 1: constants.} By Lemma~\ref{lem:flat}(a), a trajectory with $uw>0$ satisfies
\begin{equation}\label{eq:rho}
\dot\rho=\rho^{k-1}D(\theta),\qquad\dot a=-\nabla_AF(\theta),\qquad\text{where }D:=u^kG_n,
\end{equation}
and $D(\theta^*)=k\,U_n(\theta^*;u)=2\kappa_0$. Fix $r_1\in(0,r_0]$ and $C_1$ such that, on the ball $\cB=\{|\theta-\theta^*|\le r_1\}$: $\kappa_0\le D\le3\kappa_0$; $\nabla_A^2F\succeq\frac\mu2I$; $D$ and $\nabla F$ are $C_1$-Lipschitz; $|\partial_{n'}F|\le C_1|w_{n'}|^{k-1}$ for every $n'$; and Lemma~\ref{lem:flat}(b) holds with $C=C_1$. Lemma~\ref{lem:active} applies to the trajectories in $\cB$, because $|a-\theta^*|\le r_1\le r_0$ and $\zeta_A\le r_1^k\le r_0$ there. Given $K$, let $K'=K+(1+2^k)C_0$ and $C_3=C_1(K'+1+2^k)$, and choose $\rho_1,\eps_1\in(0,1]$ with
\begin{equation}\label{eq:constants}
K'\rho_1^k+3\rho_1<r_1,\qquad C_1C_3\,\rho_1^k\le\frac{\mu\kappa_0}{12},\qquad\frac{C_1}{\kappa_0}\,\eps_1^{k-2}\le1-2^{2-k}.
\end{equation}

\emph{Step 2: a priori bounds.} In the situation of (b), let $t_1\le\infty$ be the first time at which $\rho=\rho_1$. We claim that, for $t\in[t_0,t_1]$,
\begin{equation}\label{eq:apriori}
\dot\rho\ge\kappa_0\rho^{k-1},\qquad t-t_0<\frac{\rho_0^{2-k}}{(k-2)\kappa_0},\qquad|z|\le2\eta,\qquad|a-\theta^*|\le K'\rho^k .
\end{equation}
Let $[t_0,t^*]$, with $t^*\le t_1$, be the largest interval on which $\theta\in\cB$. On it, $\dot\rho\ge\kappa_0\rho^{k-1}$ by \eqref{eq:rho}. So $\rho$ increases and $\frac{d}{dt}\rho^{2-k}\le-(k-2)\kappa_0$, which bounds $t-t_0$ because $\rho^{2-k}>0$. Next, $|\dot w_{n'}|\le C_1|w_{n'}|^{k-1}$ for $n'\notin A'$, so $\frac{d}{dt}|z|\le C_1|z|^{k-1}$ and $\frac{d}{dt}|z|^{2-k}\ge-(k-2)C_1$. Hence
\[
|z(t)|^{2-k}>\eta^{2-k}-\frac{C_1}{\kappa_0}\,\rho_0^{2-k}\ge\eta^{2-k}\Big(1-\frac{C_1}{\kappa_0}\,\eps_1^{k-2}\Big)\ge(2\eta)^{2-k}
\]
by \eqref{eq:constants}; if $\eta=0$, then $z\equiv0$ because $S_{A'}$ is invariant. Finally, $\zeta_A(\theta)\le\rho^k+|z|^k$ and $\rho$ increases, so Lemma~\ref{lem:active} gives
\[
|a(t)-\theta^*|\le|a(t_0)-\theta^*|+C_0\big(\rho(t)^k+2^k\eta^k\big)\le K'\rho(t)^k,
\]
because $\eta\le\rho_0\le\rho(t)$. By \eqref{eq:constants}, these bounds give $|\theta-\theta^*|\le K'\rho^k+\rho+2\eta<r_1$ on $[t_0,t^*]$. Hence $t^*=t_1$, and $t_1<\infty$.

\emph{Step 3: comparison at equal $\rho$.} For $0<\rho\le\rho_1$, let $\cR(\rho)$ be the set of the points $\theta=a+u\rho e_n+z$ with $|a-\theta^*|\le K'\rho^k$ and $|z|\le2\eps_1\rho$. It is contained in $\cB$ by \eqref{eq:constants}. By \eqref{eq:apriori}, the trajectory of (b) lies in $\cR(\rho(t))$, and since $\rho$ increases we may use it as the time variable:
\[
\frac{da}{d\rho}=\rho^{1-k}X(\theta),\qquad X:=-\frac{\nabla_AF}{D}.
\]
On $\cR(\rho)$, Lemma~\ref{lem:flat}(b) and $\nabla F_A(\theta^*)=0$ give $|\nabla_AF(\theta)|\le C_1(\rho^k+|z|^k)+C_1|a-\theta^*|\le C_3\rho^k$. Hence, for $v\in S_A$, by \eqref{eq:constants}, with $\lambda:=\mu/(12\kappa_0)$,
\[
\ip{v}{D_aX(\theta)\,v}=-\frac{\ip{v}{\nabla_A^2F\,v}}{D}+\frac{\ip{v}{\nabla_AF}\ip{\nabla_AD}{v}}{D^2}\le\Big(-\frac{\mu}{6\kappa_0}+\frac{C_1C_3\rho^k}{\kappa_0^2}\Big)|v|^2\le-\lambda|v|^2,
\]
and, by Lemma~\ref{lem:flat}(b) for $A'$,
\[
|X(\theta)-X(\pi_{A'}\theta)|\le\frac{C_1|z|^k}{\kappa_0}+\frac{C_3\rho^k\,C_1|z|}{\kappa_0^2}\le C_4|z| .
\]
Now let $\theta_1(\rho)\in\cR(\rho)$ and $\theta_2(\rho)\in\cR(\rho)\cap S_{A'}$, $\rho\in[\rho_a,\rho_b]$, be two trajectories parametrized by $\rho$, and let $\delta=a_1-a_2$. Splitting $X(\theta_1)-X(\theta_2)=[X(\theta_1)-X(\pi_{A'}\theta_1)]+[X(\pi_{A'}\theta_1)-X(\theta_2)]$, and noting that the constraint on $a$ in $\cR(\rho)$ is convex, we get
\[
\frac{d}{d\rho}\frac{|\delta|^2}2=\rho^{1-k}\ip{\delta}{X(\theta_1)-X(\theta_2)}\le\rho^{1-k}|\delta|\big(C_4|z_1|-\lambda|\delta|\big).
\]
So $|\delta|$ decreases whenever it exceeds $C_4|z_1|/\lambda$, and
\begin{equation}\label{eq:comparison}
|\delta(\rho)|\le\max\Big\{|\delta(\rho_a)|,\ \frac{C_4}{\lambda}\sup|z_1|\Big\}\qquad\text{for }\rho\in[\rho_a,\rho_b].
\end{equation}

\emph{Step 4: the escape orbit, proof of (a).} For $\rho_0<\rho_1$, the trajectory through $\theta^*+u\rho_0e_n$ satisfies the assumptions of (b) with $\eta=0$. By Step~2 and Lemma~\ref{lem:active}, it stays in $S_{A'}$ with $|a-\theta^*|\le C_0\rho^k$. Let $a^{\rho_0}(\rho)$, $\rho\in[\rho_0,\rho_1]$, be its parametrization by $\rho$. By \eqref{eq:comparison}, for $\rho_0'<\rho_0\le\rho$,
\[
|a^{\rho_0'}(\rho)-a^{\rho_0}(\rho)|\le|a^{\rho_0'}(\rho_0)-\theta^*|\le C_0\rho_0^k .
\]
Hence, as $\rho_0\to0$, $a^{\rho_0}$ converges uniformly on every $[\rho',\rho_1]$ to a solution $a^\gamma$ of $\frac{da}{d\rho}=\rho^{1-k}X(a+u\rho e_n)$ on $(0,\rho_1]$, with $|a^\gamma(\rho)-\theta^*|\le C_0\rho^k$. Let $t(\rho)=-\int_\rho^{\rho_1}r^{1-k}/D\,dr$, with $D$ evaluated at $a^\gamma(r)+ure_n$. Since $D\le3\kappa_0$, $t(\rho)\to-\infty$ as $\rho\to0$. The curve $\gamma(t(\rho))=a^\gamma(\rho)+u\rho e_n$ solves \eqref{eq:rho}, and the other components of $\nabla F$ vanish on $S_{A'}$. So $\gamma$ is a gradient-flow trajectory with the stated properties.

Let $\tilde\gamma$ be as in (a), and $\tilde\rho=|\tilde\gamma_n|$. For $t$ near $-\infty$, $\tilde\gamma$ lies in $\cB$, so $\tilde\rho$ increases, and Lemma~\ref{lem:active} on $[t_0,t]$ with $t_0\to-\infty$ gives $|\pi_A\tilde\gamma-\theta^*|\le C_0\tilde\rho^k$. Thus \eqref{eq:comparison} applies to $\tilde\gamma$ and $\gamma$, parametrized by $\rho$ as $\tilde a(\rho)$ and $a^\gamma(\rho)$: for small $\rho$ and every $\rho_a<\rho$, $|\tilde a(\rho)-a^\gamma(\rho)|\le2C_0\rho_a^k$. Hence $\tilde a=a^\gamma$. Since $t$ is determined by $\rho$ up to an additive constant, $\tilde\gamma$ is a time translate of $\gamma$.

\emph{Step 5: proof of (b).} The first assertions are in \eqref{eq:apriori}. By \eqref{eq:comparison}, applied to $\theta$ and $\gamma$ from $\rho_a=\rho_0$,
\[
|\theta(t)-\gamma(t')|\le|a-a^\gamma|(\rho(t))+|z(t)|\le\max\Big\{(K+C_0)\rho_0^k,\ \frac{2C_4}{\lambda}\,\eta\Big\}+2\eta .\qedhere
\]
\end{proof}

\subsection*{Proof of Theorem~\ref{thm:main}}

Fix $\xi$, $M$ and $S$. For $m\le M$, let $\mu_m$, $r^{(m)}_0$ and $C^{(m)}_0$ be the constants of Lemma~\ref{lem:active} at $(A_m,\theta^{(m)})$. For $m<M$, apply Lemma~\ref{lem:escape} to $(A_m,\theta^{(m)},n_{m+1},u_{n_{m+1}})$ with $K=2^{k+1}LC^{(m)}_0$. It gives constants $\rho_1,\eps_1,C$ and $\kappa_0$, and the trajectory $\gamma=\gamma_{m+1}$, normalized by $|\gamma_{n_{m+1}}(0)|=\rho_1$. All the times $\tau$ below satisfy $\alpha^{k-2}\tau\to s$ for some $s$, and the limit objects are uniformly continuous. So statements about $\theta^\alpha(\tau+\alpha^{2-k}(s'-s))$ for $s'$ in a compact interval transfer to $\theta^\alpha(\alpha^{2-k}s')$ on slightly smaller intervals. We use this without comment. We prove by induction on $m\le M$:

\smallskip\noindent$(\mathcal C_m)$\quad There are times $\tau^\alpha_m$ with $\alpha^{k-2}\tau^\alpha_m\to s_m$, $|\pi_{A_m}\theta^\alpha(\tau^\alpha_m)-\theta^{(m)}|\le r^{(m)}_0/2$, and $\alpha^{-1}w^\alpha_n(\tau^\alpha_m)\to\Xi_n(s_m)$ for all $n\notin A_m$.

\smallskip\noindent$(\mathcal C_0)$ holds with $\tau^\alpha_0=0$.

\emph{Plateau.} Assume $(\mathcal C_m)$, and let $s_+\in(s_m,s_{m+1})$, with $s_+>S$ if $m=M$. Apply Lemma~\ref{lem:dormant} on $[s_m,s_+]$, with $R$ the maximum of the $|\Xi_n|$, $n\notin A_m$, on this interval. With $t(s)=\tau^\alpha_m+\alpha^{2-k}(s-s_m)$, it gives:
\begin{enumerate}[label=(P\arabic*),leftmargin=3em,nosep]
\item $\alpha^{-1}w^\alpha_n(t(s))\to\Xi_n(s)$ uniformly on $[s_m,s_+]$, for $n\notin A_m$;
\item for every $\delta>0$, $|\theta^\alpha(t(s))-\theta^{(m)}|\le e^{-\mu_m\alpha^{2-k}\delta}\,r^{(m)}_0/2+O(\alpha)\to0$ uniformly on $[s_m+\delta,s_+]$. For $m=0$ this holds on all of $[0,s_+]$, since $\pi_\emptyset\theta^\alpha=0$.
\end{enumerate}

\emph{Jump.} Let $m<M$, assume $(\mathcal C_m)$, and write $n_*=n_{m+1}$. Let $R'$ be the maximum of the $|\Xi_n|$, $n\notin A_{m+1}$, on $[s_m,s_{m+1}]$, and let $R_\delta=|\Xi_{n_*}(s_{m+1}-\delta)|$. By \eqref{eq:dormant}, $|\Xi_{n_*}|^{2-k}$ decreases linearly to $0$ on $[s_m,s_{m+1}]$, so $|\Xi_{n_*}|$ increases to $\infty$. Fix $\delta>0$ so small that
\[
R_\delta\ge\max\{3,R'\}\qquad\text{and}\qquad\sqrt L\,R'+1\le\eps_1(R_\delta-1).
\]
Apply the plateau step with $s_+=s_{m+1}-\delta$, so that $R=R_\delta$, and let $t_0=t(s_{m+1}-\delta)$. For small $\alpha$, by (P1) and Lemma~\ref{lem:dormant}(b):
\begin{itemize}[nosep]
\item $\rho_0:=|w^\alpha_{n_*}(t_0)|=\alpha(R_\delta+o(1))<\rho_1$, and $|w^\alpha_{n_*}|\le\alpha(R_\delta+1)<\rho_1$ on $[\tau^\alpha_m,t_0]$; the sign of $w^\alpha_{n_*}$ is $u_{n_*}$ at all times;
\item $\eta:=|z(t_0)|\le\alpha(\sqrt L\,R'+o(1))\le\eps_1\rho_0$;
\item $|a(t_0)-\theta^{(m)}|\le e^{-\mu_m\alpha^{2-k}(s_{m+1}-\delta-s_m)}\,r^{(m)}_0/2+C^{(m)}_0L(R_\delta+1)^k\alpha^k\le K\rho_0^k$. Indeed, the first term is $o(\alpha^k)$, and $(R_\delta+1)\alpha\le2(R_\delta-1)\alpha\le2\rho_0$.
\end{itemize}
So Lemma~\ref{lem:escape}(b) applies. Let $t^\alpha_{m+1}:=t_1$ be the first time after $\tau^\alpha_m$ at which $|w^\alpha_{n_*}|=\rho_1$; this definition does not involve $\delta$. We obtain:
\begin{enumerate}[label=(J\arabic*),leftmargin=3em,nosep]
\item $\alpha^{k-2}(t_1-t_0)\le(\alpha^{-1}\rho_0)^{2-k}/((k-2)\kappa_0)=(R_\delta+o(1))^{2-k}/((k-2)\kappa_0)$;
\item $|\theta^\alpha(t_1)-\gamma(0)|\le C(\rho_0^k+\eta)\to0$;
\item on $[t_0,t_1]$, $|z|\le2\eta$. So, for $n\notin A_{m+1}$, $|\frac{d}{ds}(\alpha^{-1}w^\alpha_n)|=\alpha^{1-k}|\partial_nF(\theta^\alpha)|\le C(R')$ in the time $s=\alpha^{k-2}t$, by Lemma~\ref{lem:flat}(a). By (J1), these rescaled weights move by at most $C(R')R_\delta^{2-k}+o(1)$ during $[t_0,t_1]$.
\end{enumerate}
By (J2) and the continuous dependence of solutions on their initial value, $\theta^\alpha(t_1+\cdot)\to\gamma$ uniformly on compact subsets of $\R$. By \ND3, there is $T>0$ with $|\gamma(t)-\theta^{(m+1)}|\le r^{(m+1)}_0/4$ for $t\ge T$. Set $\tau^\alpha_{m+1}=t_1+T$, so that $|\pi_{A_{m+1}}\theta^\alpha(\tau^\alpha_{m+1})-\theta^{(m+1)}|\le r^{(m+1)}_0/4+o(1)$. On $[t_1,t_1+T]$, the weights $w^\alpha_n$, $n\notin A_{m+1}$, stay $O(\alpha)$, hence they move by $O(\alpha^{k-1})$, again by Lemma~\ref{lem:flat}(a). Together with (P1), (J1) and (J3), this gives, for $n\notin A_{m+1}$,
\begin{gather*}
s_{m+1}-\delta\le\liminf_{\alpha\to0}\alpha^{k-2}\tau^\alpha_{m+1}\le\limsup_{\alpha\to0}\alpha^{k-2}\tau^\alpha_{m+1}\le s_{m+1}-\delta+\frac{R_\delta^{2-k}}{(k-2)\kappa_0},\\
\limsup_{\alpha\to0}\big|\alpha^{-1}w^\alpha_n(\tau^\alpha_{m+1})-\Xi_n(s_{m+1})\big|\le C(R')R_\delta^{2-k}+\sup_{[s_{m+1}-\delta,\,s_{m+1}]}|\Xi_n-\Xi_n(s_{m+1})| .
\end{gather*}
Since $\tau^\alpha_{m+1}$ does not depend on $\delta$, and $R_\delta\to\infty$ as $\delta\to0$, this proves $(\mathcal C_{m+1})$.

\emph{Conclusion.} By induction, the plateau and jump steps hold at all stages, and they cover $[0,\alpha^{2-k}S]$. The trajectory stays bounded there, so $\theta^\alpha$ is defined on this interval. Statement (i) is (P2), and (ii) follows from (P1), together with (J3) and the last display for the dormant neurons during a jump. The loss converges because $F$ is continuous, and $F(\theta^{(m)})>F(\theta^{(m+1)})$ because $F$ decreases strictly along $\gamma_{m+1}$. Each $\theta^{(m)}$ is a critical point of $F$, and a saddle point for $m<M$, by the consequences of Lemma~\ref{lem:flat}, since $U^{(m)}_{n_{m+1}}(u_{n_{m+1}})>0$. The local uniform convergence in (iii) was shown in the jump step, with $t^\alpha_{m+1}=t_1$.

For the Hausdorff convergence at $s_{m+1}$, let $0<\eps\le r^{(m+1)}_0/4$, and choose $T',T>0$ with $\gamma((-\infty,-T'])\subseteq B(\theta^{(m)},\eps)$ and $\gamma([T,\infty))\subseteq B(\theta^{(m+1)},\eps)$. Run the jump step with this $T$, and with some $\delta'<\delta$ in place of $\delta$. The curve $\Gamma^\alpha=\{\theta^\alpha(\alpha^{2-k}s):|s-s_{m+1}|\le\delta\}$ consists of four pieces. Before $t_0$, it is $o(1)$-close to $\theta^{(m)}$, by (P2). On $[t_0,t_1]$, it is $o(1)$-close to $\gamma((-\infty,0])$, by Lemma~\ref{lem:escape}(b). On $[t_1,t_1+T]$, it is $o(1)$-close to $\gamma([0,T])$. Afterwards, it is $(\eps+o(1))$-close to $\theta^{(m+1)}$, by Lemma~\ref{lem:active} and (P1) at stage $m+1$. Conversely, every point of $\gamma(\R)\cup\{\theta^{(m)},\theta^{(m+1)}\}$ is $(\eps+o(1))$-close to $\Gamma^\alpha$: the points of $\gamma([-T',T])$ by the local uniform convergence, and the others by (P2) at $s_{m+1}\pm\delta$. Hence $\limsup_{\alpha\to0}d_H\big(\Gamma^\alpha,\gamma(\R)\cup\{\theta^{(m)},\theta^{(m+1)}\}\big)\le\eps$ for every small $\eps>0$.\qed

\section{Remarks}

\begin{remark}[Random initialization]
Theorem~B of the full note simplifies in the same way. Let $\cK$ be the set of the configurations $(A,\theta)$ that are reached from $(\emptyset,0)$ along escape orbits converging to non-degenerate local minima, as in the full note. It is finite, since a configuration has at most $2L$ successors, and condition (G2) is void. Assume (G1): for every $(A,\theta)\in\cK$, every $n\notin A$ and every $u\in\{\pm1\}$ with $U_n(\theta;u)>0$, the escape orbit of $\theta$ along $(n,u)$ is defined for all $t\ge0$ and converges to a non-degenerate local minimum of $F_{A\cup\{n\}}$. Then every stage is non-degenerate, unless $\xi$ lies in a finite union of hypersurfaces. In particular, this holds almost surely if $\xi$ has a density. Indeed, \ND3 follows from (G1). For \ND1, fix the signs $u_n$ and the first activations $n_1,\dots,n_m$. They determine $(A_j,\theta^{(j)})_{j\le m}$, hence the speeds of the clocks. By Remark~\ref{rem:race}, the times $s_1,\dots,s_m$ are then linear functions of $y=(|\xi_n|^{2-k})_n$, and the clock of a dormant neuron $n$ at $s_m$ is $y_n$ minus a linear function of $y_{n_1},\dots,y_{n_m}$. So a tie between two dormant neurons at stage $m$ is a non-trivial linear equation in $y$, and there are finitely many such equations. Similarly, a fixed $s>0$ is a jump time only if $y$ solves one of finitely many non-trivial affine equations. So $\theta^\alpha(\alpha^{2-k}s)\to\bar\theta(s)$ almost surely, for every $s\ge0$.
\end{remark}

\begin{remark}[Biases]
With biases, $h_n=h_{n-1}+\sigma(w_nh_{n-1}+b_n)$, the incoming weights of a neuron form a vector $(w_n,b_n)\in\R^2$, the loss is flat along $\{(w_n,b_n)=0\}$, and the utility is a binary form of degree $k$. The direction of a neuron is then a point of the circle, and it moves, by gradient ascent of the utility. So the dormant flow is no longer explicit, \ND2 is a genuine condition, and the escape lemma has the direction as an additional variable, which is contracted only at rate $O(1)$ in the time $\log\rho$. This is the general proof with $p=2$, for which we refer to the full note.
\end{remark}

\end{document}
